\pagestyle{empty}
\tight
\begin{tblr}{width=\textwidth,colspec={|Q[c,m,1.9cm]|X[l,m]|},hlines,vlines,hspan=minimal,row{1}={bg=headblue},rowsep=1pt,colsep=3pt}
\SetCell{c,m}\hfil\logo\hfil & \textbf{POLITEKNIK NEGERI MALANG}\par\textbf{JURUSAN TEKNOLOGI INFORMASI}\par\textbf{PROGRAM STUDI: D4 TEKNIK INFORMATIKA} \\
\SetCell[c=2]{c} \textbf{RENCANA PEMBELAJARAN SEMESTER (RPS)} & \\
\end{tblr}
\begin{longtblr}{width=\textwidth,colspec={|Q[l,m,1.9cm]|X[0.85,l,m]|X[1.45,l,m]|X[0.75,l,m]|X[1.15,l,m]|},hlines,vlines,hspan=minimal,rowsep=1pt,colsep=2pt,row{1,3}={bg=headgray,font=\bfseries}}
MATA KULIAH & KODE & BOBOT (SKS)/jam & SEMESTER & TGL. PENYUSUNAN \\
Analisis dan Desain Berorientasi Objek & RTI254002 & 2 SKS/4 jam & 4 & 1 Juli 2025 \\
OTORISASI & Dosen Pengembang RPS & Koordinator RMK & \SetCell[c=2]{l} Ka PRODI &  \\
 & Agung Nugroho Pramudhita, S.T., M.T. &  & \SetCell[c=2]{l}{\vspace{8mm}\par Dr. Ely Setyo Astuti, S.T., M.T.} &  \\
\SetCell[r=12]{l} Capaian Pembelajaran (CP) & \SetCell[c=4]{l,bg=headgray,font=\bfseries} Capaian Pembelajaran Lulusan Program Studi (CPL-Prodi) &  &  &  \\
 & CPL02 & \SetCell[c=3]{l} Mampu menerapkan prinsip rekayasa sistem dan perangkat lunak untuk menghasilkan solusi aplikasi multi-platform sesuai kebutuhan. &  &  \\
 & CPL06 & \SetCell[c=3]{l} Mampu merancang, mengimplementasikan, menguji, melakukan deployment, memelihara, serta menjamin mutu perangkat lunak yang menjawab permasalahan dan memenuhi kebutuhan pengguna. &  &  \\
 & CPL10 & \SetCell[c=3]{l} Mampu menganalisis persoalan kompleks dengan wawasan multidisiplin untuk menghasilkan solusi di bidang informatika dan ilmu komputer. &  &  \\
 & \SetCell[c=4]{l,bg=headgray,font=\bfseries} Capaian Pembelajaran mata kuliah (CPL-MK) &  &  &  \\
 & CPMK0211 & \SetCell[c=3]{l} Mahasiswa mampu menerapkan metodologi analisis dan desain berorientasi objek untuk menghasilkan spesifikasi teknis komprehensif dari sistem perangkat lunak. &  &  \\
 & CPMK0602 & \SetCell[c=3]{l} Mahasiswa mampu merancang komponen perangkat lunak menggunakan software design pattern yang tepat sesuai kebutuhan sistem. &  &  \\
 & CPMK1008 & \SetCell[c=3]{l} Mahasiswa mampu mengintegrasikan hasil analisis dan desain sistem berorientasi objek ke dalam spesifikasi pengembangan yang dapat diimplementasikan. &  &  \\
 & \SetCell[c=4]{l,bg=headgray,font=\bfseries} Kemampuan akhir tiap tahapan belajar (Sub-CPMK) &  &  &  \\
 & SCPMK0211-02901 & \SetCell[c=3]{l} Mahasiswa mampu menganalisis kebutuhan sistem dan memodelkan domain menggunakan diagram UML (use case, class, sequence, activity). &  &  \\
 & SCPMK0602-02902 & \SetCell[c=3]{l} Mahasiswa mampu merancang arsitektur perangkat lunak berorientasi objek menggunakan design pattern (GoF) dan prinsip SOLID. &  &  \\
 & SCPMK1008-02903 & \SetCell[c=3]{l} Mahasiswa mampu mengintegrasikan artefak ADBO (SRS, model domain, diagram arsitektur) menjadi spesifikasi teknis yang siap diimplementasikan. &  &  \\
\SetCell[r=2]{l} Penilaian & \SetCell[c=4]{l} *Nilai CPMK didapat dari agregasi nilai SUB-CPMK yang dibebankan &  &  &  \\
 & \SetCell[c=4]{l}{\tiny\begin{tblr}{width=\linewidth,colspec={|Q[c,m,0.1\linewidth]|Q[c,m,0.16\linewidth]|X[l,m]|X[l,m]|X[l,m]|X[l,m]|X[l,m]|X[l,m]|Q[c,m,0.08\linewidth]|},hlines,vlines,hspan=minimal,rowsep=1pt,colsep=0.7pt,row{1,2}={bg=headgray,font=\bfseries}}
Kode CPMK & Kode SCPMK & \SetCell[c=6]{c} Bobot per Bentuk Penilaian & & & & & & Total Bobot \\
 & & Tugas & Kuis 1 & UTS & Kuis 2 & PBL & UAS & \\
CPMK0211 & SCPMK0211-02901 & 12\%: artefak analisis UML & 5\%: UML dan kasus & 8\%: UTS artefak analisis & & & & 25\% \\
CPMK0602 & SCPMK0602-02902 & 10\%: tugas design pattern & & 5\%: UTS design pattern & 5\%: kuis design pattern & 10\%: PBL desain terpadu & & 30\% \\
CPMK1008 & SCPMK1008-02903 & 8\%: artefak spesifikasi teknis & & & & 12\%: PBL desain terpadu & 25\%: review dokumen desain & 45\% \\
\SetCell[c=8]{r}\textbf{Total Per Penilaian} & & & & & & & & \textbf{100\%} \\
\end{tblr}} &  &  &  \\
Diskripsi Singkat Mata Kuliah & \SetCell[c=4]{l} Mata kuliah ini membangun kemampuan menganalisis kebutuhan sistem dan merancang solusi perangkat lunak berorientasi objek menggunakan UML, design pattern GoF, dan prinsip SOLID untuk menghasilkan spesifikasi teknis yang implementable. &  &  &  \\
Bahan Kajian / Materi Pembelajaran & \SetCell[c=4]{l}\begin{enumerate}\item Paradigma OO, requirement engineering, dan use case modeling\item Domain modeling, class diagram, sequence, state, dan activity diagram\item Prinsip SOLID dan pengenalan design patterns\item Creational, structural, dan behavioral design patterns (GoF)\item Arsitektur MVC, component diagram, dan deployment diagram\end{enumerate} &  &  &  \\
\SetCell[r=2]{l} Pustaka & \SetCell[c=4]{l} \textbf{Utama:}\begin{enumerate}\item Larman, C. 2005. Applying UML and Patterns. Prentice Hall.\item Gamma, E. 1994. Design Patterns: Elements of Reusable OO Software. Addison-Wesley.\item Fowler, M. 2004. UML Distilled. Addison-Wesley.\end{enumerate} &  &  &  \\
 & \SetCell[c=4]{l} \textbf{Pendukung:} Materi LMS, dokumentasi PlantUML/StarUML, dan jobsheet analisis desain. &  &  &  \\
Media Pembelajaran & \SetCell[c=2]{l} \textbf{Software:} draw.io/StarUML/PlantUML, office suite, LMS. &  & \SetCell[c=2]{l} \textbf{Hardware:} Komputer/laptop, proyektor. &  \\
Nama Dosen Pengampu & \SetCell[c=4]{l} Tim Pengajar Program Studi D4 TI &  &  &  \\
Matakuliah Syarat & \SetCell[c=4]{l} RTI252005 Rekayasa Perangkat Lunak &  &  &  \\
\end{longtblr}
\begin{longtblr}{width=\textwidth,colspec={|Q[c,m,0.06\textwidth]|X[1.35,l,m]|X[1.08,l,m]|X[1.16,l,m]|X[0.7,l,m]|X[1.24,l,m]|X[1.06,l,m]|X[0.98,l,m]|Q[c,m,0.05\textwidth]|},hlines,vlines,hspan=minimal,rowhead=2,row{1,2}={bg=headgreen,font=\bfseries},rowsep=1pt,colsep=1.5pt}
Mgg.\par Ke & Kemampuan Akhir Yang Direncanakan (Sub-CP-MK) & Bahan kajian (Materi Pembelajaran) & Bentuk dan Metode Pembelajaran & Estimasi Waktu & Pengalaman Belajar Mahasiswa & Kriteria \& Bentuk Penilaian & Indikator Penilaian & Bobot Penilaian (\%) \\
(1) & (2) & (3) & (4) & (5) & (6) & (7) & (8) & (9) \\
1 & SCPMK0211-02901 & Pengantar ADBO: paradigma OO dan proses analisis desain. & Modalitas: Blended Learning. Bentuk: Luring/Daring. Metode: Case Method, analisis kasus, studi literatur, dan presentasi. & 1 x 2 x 50' tatap muka; 1 x 2 x 50' tugas/praktik mandiri. & Mahasiswa mempelajari paradigma OO dan mengerjakan tugas analisis kebutuhan awal dari studi kasus. & Bentuk penilaian: Tugas Artefak. Mengacu rubrik RTM. & Ketepatan konsep; kualitas artefak; validasi dan dokumentasi. & 2 \\
2 & SCPMK0211-02901 & Requirement engineering dan use case modeling. & Modalitas: Blended Learning. Bentuk: Luring/Daring. Metode: Case Method, analisis kasus, studi literatur, dan presentasi. & 1 x 2 x 50' tatap muka; 1 x 2 x 50' tugas/praktik mandiri. & Mahasiswa mengidentifikasi kebutuhan fungsional dan memodelkan use case dari studi kasus yang diberikan. & Bentuk penilaian: Tugas Artefak. Mengacu rubrik RTM. & Ketepatan konsep; kualitas artefak; validasi dan dokumentasi. & 2 \\
3 & SCPMK0211-02901 & Use case specification dan scenario analysis. & Modalitas: Blended Learning. Bentuk: Luring/Daring. Metode: Case Method, analisis kasus, studi literatur, dan presentasi. & 1 x 2 x 50' tatap muka; 1 x 2 x 50' tugas/praktik mandiri. & Mahasiswa menulis spesifikasi use case lengkap dengan skenario utama dan alternatif. & Bentuk penilaian: Tugas Artefak. Mengacu rubrik RTM. & Ketepatan konsep; kualitas artefak; validasi dan dokumentasi. & 2 \\
4 & SCPMK0211-02901 & Domain modeling dan class diagram. & Modalitas: Blended Learning. Bentuk: Luring/Daring. Metode: Case Method, analisis kasus, studi literatur, dan presentasi. & 1 x 2 x 50' tatap muka; 1 x 2 x 50' tugas/praktik mandiri. & Mahasiswa membangun model domain dan class diagram awal dari studi kasus yang dianalisis. & Bentuk penilaian: Tugas Artefak. Mengacu rubrik RTM. & Ketepatan konsep; kualitas artefak; validasi dan dokumentasi. & 2 \\
5 & SCPMK0211-02901 & Kuis 1: use case dan class diagram. & Modalitas: Blended Learning. Bentuk: Luring/Daring. Metode: Case Method, analisis kasus, studi literatur, dan presentasi. & 1 x 2 x 50' tatap muka; 1 x 2 x 50' tugas/praktik mandiri. & Mahasiswa mengerjakan kuis untuk mengukur pemahaman use case modeling dan class diagram. & Bentuk penilaian: Kuis 1. Mengacu rubrik RTM. & Ketepatan konsep; kualitas artefak; validasi dan dokumentasi. & 5 \\
6 & SCPMK0211-02901 & Sequence diagram dan collaboration diagram. & Modalitas: Blended Learning. Bentuk: Luring/Daring. Metode: Case Method, analisis kasus, studi literatur, dan presentasi. & 1 x 2 x 50' tatap muka; 1 x 2 x 50' tugas/praktik mandiri. & Mahasiswa merancang sequence diagram dan collaboration diagram untuk use case yang telah dianalisis. & Bentuk penilaian: Tugas Artefak. Mengacu rubrik RTM. & Ketepatan konsep; kualitas artefak; validasi dan dokumentasi. & 2 \\
7 & SCPMK0211-02901 & State diagram dan activity diagram. & Modalitas: Blended Learning. Bentuk: Luring/Daring. Metode: Case Method, analisis kasus, studi literatur, dan presentasi. & 1 x 2 x 50' tatap muka; 1 x 2 x 50' tugas/praktik mandiri. & Mahasiswa merancang state diagram dan activity diagram untuk alur proses pada studi kasus. & Bentuk penilaian: Tugas Artefak. Mengacu rubrik RTM. & Ketepatan konsep; kualitas artefak; validasi dan dokumentasi. & 2 \\
8 & SCPMK0211-02901, SCPMK0602-02902 & UTS: artefak analisis ADBO. & Modalitas: Blended Learning. Bentuk: Luring/Daring. Metode: Case Method, analisis kasus, studi literatur, dan presentasi. & 1 x 2 x 50' tatap muka; 1 x 2 x 50' tugas/praktik mandiri. & Mahasiswa mengerjakan ujian tengah semester dengan menghasilkan artefak analisis ADBO dari kasus baru yang diberikan. & Bentuk penilaian: UTS. Mengacu rubrik RTM. & Ketepatan konsep; kualitas artefak; validasi dan dokumentasi. & 13 \\
9 & SCPMK0602-02902 & Prinsip SOLID dan pengenalan design patterns. & Modalitas: Blended Learning. Bentuk: Luring/Daring. Metode: Case Method, analisis kasus, studi literatur, dan presentasi. & 1 x 2 x 50' tatap muka; 1 x 2 x 50' tugas/praktik mandiri. & Mahasiswa menganalisis prinsip SOLID pada kode yang ada dan mengidentifikasi pelanggaran serta solusinya. & Bentuk penilaian: Tugas Artefak. Mengacu rubrik RTM. & Ketepatan konsep; kualitas artefak; validasi dan dokumentasi. & 2.5 \\
10 & SCPMK0602-02902 & Creational patterns: Singleton, Factory, Builder. & Modalitas: Blended Learning. Bentuk: Luring/Daring. Metode: Case Method, analisis kasus, studi literatur, dan presentasi. & 1 x 2 x 50' tatap muka; 1 x 2 x 50' tugas/praktik mandiri. & Mahasiswa menerapkan creational patterns pada studi kasus yang relevan dan mendokumentasikan hasilnya. & Bentuk penilaian: Tugas Artefak. Mengacu rubrik RTM. & Ketepatan konsep; kualitas artefak; validasi dan dokumentasi. & 2.5 \\
11 & SCPMK0602-02902 & Structural patterns: Adapter, Decorator, Facade. & Modalitas: Blended Learning. Bentuk: Luring/Daring. Metode: Case Method, analisis kasus, studi literatur, dan presentasi. & 1 x 2 x 50' tatap muka; 1 x 2 x 50' tugas/praktik mandiri. & Mahasiswa menerapkan structural patterns pada studi kasus yang relevan dan mendokumentasikan hasilnya. & Bentuk penilaian: Tugas Artefak. Mengacu rubrik RTM. & Ketepatan konsep; kualitas artefak; validasi dan dokumentasi. & 2.5 \\
12 & SCPMK0602-02902 & Behavioral patterns: Observer, Strategy, Command. & Modalitas: Blended Learning. Bentuk: Luring/Daring. Metode: Case Method, analisis kasus, studi literatur, dan presentasi. & 1 x 2 x 50' tatap muka; 1 x 2 x 50' tugas/praktik mandiri. & Mahasiswa menerapkan behavioral patterns pada studi kasus yang relevan dan mendokumentasikan hasilnya. & Bentuk penilaian: Tugas Artefak. Mengacu rubrik RTM. & Ketepatan konsep; kualitas artefak; validasi dan dokumentasi. & 2.5 \\
13 & SCPMK0602-02902 & Kuis 2: design patterns. & Modalitas: Blended Learning. Bentuk: Luring/Daring. Metode: Case Method, analisis kasus, studi literatur, dan presentasi. & 1 x 2 x 50' tatap muka; 1 x 2 x 50' tugas/praktik mandiri. & Mahasiswa mengerjakan kuis untuk mengukur pemahaman design patterns GoF dan penerapannya. & Bentuk penilaian: Kuis 2. Mengacu rubrik RTM. & Ketepatan konsep; kualitas artefak; validasi dan dokumentasi. & 5 \\
14 & SCPMK1008-02903 & Arsitektur MVC dan layered architecture. & Modalitas: Blended Learning. Bentuk: Luring/Daring. Metode: Case Method, analisis kasus, studi literatur, dan presentasi. & 1 x 2 x 50' tatap muka; 1 x 2 x 50' tugas/praktik mandiri. & Mahasiswa merancang arsitektur MVC dan layered architecture untuk studi kasus dan mendokumentasikan hasilnya. & Bentuk penilaian: Tugas Artefak. Mengacu rubrik RTM. & Ketepatan konsep; kualitas artefak; validasi dan dokumentasi. & 4 \\
15 & SCPMK1008-02903 & Component diagram dan deployment diagram. & Modalitas: Blended Learning. Bentuk: Luring/Daring. Metode: Case Method, analisis kasus, studi literatur, dan presentasi. & 1 x 2 x 50' tatap muka; 1 x 2 x 50' tugas/praktik mandiri. & Mahasiswa merancang component diagram dan deployment diagram yang terintegrasi dengan artefak ADBO sebelumnya. & Bentuk penilaian: Tugas Artefak. Mengacu rubrik RTM. & Ketepatan konsep; kualitas artefak; validasi dan dokumentasi. & 4 \\
16 & SCPMK0602-02902, SCPMK1008-02903 & PBL: desain sistem ADBO terpadu. & Modalitas: Blended Learning. Bentuk: Luring/Daring. Metode: Case Method, analisis kasus, studi literatur, dan presentasi. & 1 x 2 x 50' tatap muka; 1 x 2 x 50' tugas/praktik mandiri. & Mahasiswa mengintegrasikan seluruh artefak ADBO menjadi dokumen desain sistem lengkap dan mempresentasikannya. & Bentuk penilaian: PBL Desain Sistem. Mengacu rubrik RTM. & Ketepatan konsep; kualitas artefak; validasi dan dokumentasi. & 22 \\
17 & SCPMK1008-02903 & UAS: presentasi dan review dokumen desain sistem. & Modalitas: Blended Learning. Bentuk: Luring/Daring. Metode: Case Method, analisis kasus, studi literatur, dan presentasi. & 1 x 2 x 50' tatap muka; 1 x 2 x 50' tugas/praktik mandiri. & Mahasiswa mempresentasikan dokumen desain sistem final dan mempertahankan setiap keputusan desain kepada penguji. & Bentuk penilaian: UAS. Mengacu rubrik RTM. & Ketepatan konsep; kualitas artefak; validasi dan dokumentasi. & 25 \\
\end{longtblr}
